\documentclass[10pt, aspectratio=169]{beamer}
\usetheme{Madrid}
\usecolortheme{dolphin}
\usepackage{ctex}
\usepackage{xeCJK}
\usepackage{amsmath,amssymb}
\usepackage{listings}
\usepackage{tikz}
\usepackage{xcolor}
\title{408考研零基础 --- 数据结构}
\subtitle{1-16 图的遍历}
\author{}
\date{}
\setbeamertemplate{page number in head/foot}{}
\begin{document}
\frame{\titlepage}
\begin{frame}{本节目标}
\begin{enumerate}
\item 理解 DFS 的原理与实现
\item 理解 BFS 的原理与实现
\item 掌握 DFS 和 BFS 在连通性、生成树中的应用
\end{enumerate}
\end{frame}
\begin{frame}{深度优先搜索 DFS}
\textbf{策略}：沿路径深入探索，回溯继续
\vspace{0.3em}
\textbf{递归算法}
\begin{itemize}
\item 访问 $v$，标记 visited
\item 取 $v$ 的未访问邻接顶点 $w$，递归 DFS($w$)
\item 所有邻接顶点已访问则回溯
\end{itemize}
\vspace{0.3em}
\textbf{时间复杂度}
\begin{itemize}
\item 邻接表：$O(n+e)$
\item 邻接矩阵：$O(n^2)$
\end{itemize}
空间复杂度 $O(n)$（递归栈）
\end{frame}
\begin{frame}{广度优先搜索 BFS}
\textbf{策略}：按距离递增逐层访问
\vspace{0.3em}
\textbf{队列实现}
\begin{itemize}
\item 访问 $v$，入队
\item 出队 $u$，访问 $u$ 的所有未访问邻接顶点并入队
\item 重复直到队空
\end{itemize}
\vspace{0.3em}
\textbf{时间复杂度}
\begin{itemize}
\item 邻接表：$O(n+e)$
\item 邻接矩阵：$O(n^2)$
\end{itemize}
空间复杂度 $O(n)$（队列）
\vspace{0.3em}
\textbf{应用}：无权图单源最短路径
\end{frame}
\begin{frame}{图的连通性判定}
一次 DFS/BFS 只能访问一个连通分量
\vspace{0.3em}
\textbf{遍历所有顶点}
\begin{itemize}
\item 对每个未访问顶点启动一次遍历
\item 遍历次数 $=$ 连通分量数
\end{itemize}
\vspace{0.5em}
\textbf{生成树}
\begin{itemize}
\item DFS 路径 $\to$ 深度优先生成树
\item BFS 路径 $\to$ 广度优先生成树
\item 包含全部 $n$ 个顶点，$n-1$ 条边
\end{itemize}
\end{frame}
\begin{frame}{本节总结}
\begin{enumerate}
\item DFS：递归/栈，沿路径深入，$O(n+e)$
\item BFS：队列，按层访问，$O(n+e)$
\item 一次遍历一个连通分量
\item 遍历路径构成生成树
\end{enumerate}
\vspace{1em}
\begin{center}
\textcolor{blue}{\textbf{下节预告：查找基础}}
\end{center}
\end{frame}
\end{document}
